% Introduction. Results-dependent sentences were written only after the single unseal, from results/analysis.json.
\section{Introduction}\label{sec:intro}

Stochastic perturbation of attention is a common regularizer for transformers
\citep{vaswani2017attention}. Attention-probability dropout, DropAttention \citep{lin2019dropattention},
DropKey \citep{li2023dropkey}, DropHead \citep{zhou2020drophead} and Relaxed Attention
\citep{lohrenz2021relaxed,lohrenz2023relaxed} are all specified by a nominal parameter: a drop rate $p$
or a mixing coefficient $\gamma$. For a single head, however, the forward pass is changed by the
\emph{realized} displacement of the attention row, $d=\tfrac12\lVert\tilde a-a\rVert_1$. If the values
have diameter $D_V$, the head output moves by at most $D_V d$ (Section~\ref{sec:theory}). How large $d$ is
for a given rate depends on how the row distributes its mass. For attention dropout the expected
displacement equals the rate. For targeted deletion it is the rate times the mass of the targeted keys.
In the pretrained BiomedBERT encoder \citep{gu2021domain}, the four highest-weighted keys hold
6.9\% of a row's mass at the 5th percentile of rows and 97\% at the 95th. The same top-4 deletion rate
therefore displaces some rows about fourteen times more than others, and on concentrated rows the displacement
is heavy-tailed.

This has two consequences. First, comparing families at equal nominal rates confounds the family with
the dose it delivers. Second, a nominal rate controls neither the mean nor the tail of the displacement.
We study a direct way to set the displacement instead. Given any candidate perturbed row $b$, the
capped mixture $\tilde a=(1-\alpha)a+\alpha b$ with $\alpha=\min\{1,\tau/\mathrm{TV}(a,b)\}$ displaces
every row by exactly $\min\{\mathrm{TV}(a,b),\tau\}$. With a top-$k$ deletion candidate we call this
CapDrop. It keeps the deletion events of the candidate and removes the part of the displacement above $\tau$.
The coefficient is radial norm clipping \citep{koloskova2023revisiting} of the displacement $b-a$ in the
total-variation (TV) norm. The operator is standard. What we add is the quantity it fixes during
training, and a test of whether fixing it matters.

That test needs a comparison at matched mean displacement, because otherwise any effect of the cap
could be an effect of a smaller dose. Among all mechanisms whose displacement lies pathwise between
zero and that of the candidate and has a given mean, truncation is minimal in convex order. This is
the classical optimality of stop-loss truncation \citep{ohlin1969class,kaas2008modern}, and it makes
CapDrop the least dispersed member of that class. A fixed-coefficient mixture of the same candidate
belongs to the same class. Comparing the two at the same realized mean displacement, with the same
deletion draws, isolates the tail of the displacement from its mean and from the candidate.

We run this comparison in two data-limited biomedical regimes where attention regularization is
plausibly relevant and fine-tuning is known to be fragile \citep{mosbach2021stability,dodge2020finetuning,zhang2021revisiting}.
The first (S1) is BiomedBERT fine-tuned on 450 expert-labelled PubMedQA questions per fold
\citep{jin2019pubmedqa}. The second (S2) is a small causal language model trained from scratch for
many epochs on a fixed PubMed corpus, a data-constrained setting \citep{muennighoff2023scaling}. Every
arm is calibrated to the same realized mean displacement, using a census of attention rows taken
before any outcome was observed. The study was preregistered with frozen code, data manifests and
analysis. Test labels and the language-model test stream were sealed and are used once, after every
declared run is complete, and every run is reported.

\paragraph{Contributions}
\begin{enumerate}
\item \textbf{Displacement accounting.} We give exact expressions for the expected realized displacement of
attention dropout, independent and top-$k$ key deletion and Relaxed Attention, together with the
second moments and the train/evaluation mean shift where these have closed forms
(Propositions~\ref{prop:deletion}--\ref{prop:laws}). We use realized displacement, rather than
nominal rate, as the axis on which families are compared. The individual identities are elementary,
and several are known.
\item \textbf{A TV cap for training-time perturbations.} For CapDrop and for TV-budgeted relaxation we
state the cap identities, the per-row output bound and the reading of $\tau$ as a dose whose
marginal effect is the measurable activation rate $\Pr(m>\tau)$ (Propositions~\ref{prop:cap}--\ref{thm:cx}).
The cap is closed form and adds $O(n)$ operations per row to the candidate.
\item \textbf{A preregistered, displacement-calibrated test.} A $2\times2$ design crosses the candidate
(uniform or deletion) with the mixing rule (fixed or capped). It is run alongside calibrated dropout,
DropKey and unbudgeted deletion, and tuned non-attention controls, in both regimes. Six confirmatory
contrasts use exact tests with Holm correction and equivalence margins, and claims follow a table of
outcomes fixed in advance. Declared extensions, run in a fixed order within a fixed compute budget,
cover a second seed, a fragile fine-tuning recipe, a second encoder, a second biomedical task,
R-Drop, budgeted DropKey, distilled students and a stop-gradient ablation of the cap. All of them completed within the budget, except a
single-epoch LM extension that did not fit.
\item \textbf{Evidence-use diagnostics and data hygiene for PubMedQA.} We evaluate every fine-tuned
model with the question only, with a label-differing mismatched abstract, and with the RESULTS
section deleted. We report these next to a question-only model trained under the same protocol. We also
document two unlabelled PubMedQA records that duplicate official test questions under different
PubMed identifiers.
\end{enumerate}

The results support a narrow claim. In the data-limited language model, CapDrop beat fixed-coefficient deletion by
0.045 nats/token at the same mean displacement and with the same deletion draws. It beat tuned Relaxed Attention by 0.032 nats/token
(both Holm-adjusted, on all ten seeds). We had registered the opposite expectation for this regime. A
stop-gradient ablation attributes about a fifth of the first difference to the gradient through the cap, and the
rest to the capped perturbation itself. The strongest tuned control, DropKey at twice the mean displacement, was better
than CapDrop by 0.058 nats/token, and we report this as a cost of the method. In PubMedQA fine-tuning, the
confirmatory contrasts were inconclusive, and no attention regularizer improved on the unregularized model. The only
significant difference, in favour of label smoothing, was a calibration effect that temperature scaling removed. The extensions did not
change this picture. Following the outcome table fixed before the study, we claim that the tail of the displacement
matters at a fixed mean in the language-model regime. We do not claim that capping is the active ingredient of
attention regularization in general.

Section~\ref{sec:related} positions the work. Section~\ref{sec:theory} develops the accounting,
and Section~\ref{sec:design} gives the preregistered design. Sections~\ref{sec:results}--\ref{sec:discussion}
report and discuss the results.
